\documentclass[11pt,a4paper]{article}
\usepackage[utf8]{inputenc}
\usepackage[T1]{fontenc}
\usepackage[french]{babel}
\usepackage{amsmath}
\usepackage{amssymb}
\usepackage{listings}
\usepackage{xcolor}
\usepackage{graphicx}

% Configuration des blocs de code VHDL
\lstset{
    language=VHDL,
    basicstyle=\ttfamily\footnotesize,
    keywordstyle=\color{blue}\bfseries,
    commentstyle=\color{green!60!black},
    stringstyle=\color{red},
    showstringspaces=false,
    frame=single,
    numbers=left,
    numberstyle=\tiny\color{gray},
    breaklines=true,
    tabsize=4
}

\title{TD4 Correction \\ \large SN360/SN361 : Introduction à la conception des circuits numériques}
\author{Grenoble INP - Esisar - EI3-2024/2025}
\date{}

\begin{document}

\maketitle

\section*{Exercice 1 :}

Vous allez décrire en VHDL une bascule Toggle dont la sortie $q$ change de valeur à chaque nouveau front d'horloge \texttt{clk} quand \texttt{t} est actif. Le signal de remise à zéro \texttt{reset} est asynchrone.

\begin{table}[h!]
    \centering
    \begin{tabular}{|c|c|c|c|}
        \hline
        \textbf{clk} & \textbf{t} & \textbf{reset} & $\mathbf{q'}$ \\ \hline
        Front montant & 1 & 0 & $\overline{q}$ \\ \hline
        Front montant & 0 & 0 & $q$ \\ \hline
        X & X & 1 & 0 \\ \hline
    \end{tabular}
    \caption{Table de vérité de la bascule Toggle}
\end{table}

\paragraph{1) Donner le diagramme de transitions d'une machine à états de Moore de ce circuit.}
\textit{(Le schéma SVG de la machine de Moore correspond aux états S0/S1 avec les transitions $t=0$ et $t=1$ et l'initialisation $reset=1 \rightarrow q=0$).}

\paragraph{2) Donner le code VHDL de cette machine (décrire uniquement son architecture) avec 2 processus comme vu dans le cours.}

\begin{lstlisting}[language=VHDL]
library IEEE;
use IEEE.STD_LOGIC_1164.ALL;

entity Toggle_FlipFlop_Moore is
port (
    clk: in std_logic;
    reset: in std_logic;
    t: in std_logic;
    q: out std_logic
);
end Toggle_FlipFlop_Moore;

architecture Behavioral of Toggle_FlipFlop_Moore is
    type state_type is (S0, S1);
    signal current_state, next_state: state_type := S0;
begin

    process(clk, reset)
    begin
        if reset = '1' then
            current_state <= S0;
        elsif rising_edge(clk) then
            current_state <= next_state;
        end if;
    end process;

    process(current_state, t)
    begin
        next_state <= current_state;
        case current_state is
            when S0 =>
                q <= '0';
                if t = '1' then
                    next_state <= S1;
                end if;
            when S1 =>
                q <= '1';
                if t = '1' then
                    next_state <= S0;
                end if;
            when others =>
                next_state <= S0;
        end case;
    end process;
end Behavioral;
\end{lstlisting}

\paragraph{3) Faire de même (questions 1 et 2) avec une machine de Mealy.}

\textit{(Le schéma SVG de la machine de Mealy assigne les sorties sur les transitions).}

\begin{lstlisting}[language=VHDL]
library IEEE;
use IEEE.STD_LOGIC_1164.ALL; 

entity Toggle_FlipFlop_Mealy is 
port (
    clk: in std_logic;
    reset: in std_logic;
    t: in std_logic;
    q: out std_logic
);
end Toggle_FlipFlop_Mealy; 

architecture Behavioral of Toggle_FlipFlop_Mealy is 
    type state_type is (S0, S1);
    signal current_state, next_state: state_type := S0; 
begin 
    process(clk, reset) 
    begin
        if reset = '1' then 
            current_state <= S0;
        elsif rising_edge(clk) then 
            current_state <= next_state; 
        end if;
    end process; 

    process(current_state, t) 
    begin 
        case current_state is 
            when S0 => 
                if t = '1' then 
                    next_state <= S1;
                    q <= '0'; 
                else 
                    next_state <= S0;
                    q <= '0'; 
                end if; 
            when S1 => 
                if t = '1' then 
                    next_state <= S0;
                    q <= '1'; 
                else 
                    next_state <= S1; -- Stay in S1 if t=0
                    q <= '1';
                end if;
            when others =>
                next_state <= S0;
                q <= '0';
        end case;
    end process;
end Behavioral;
\end{lstlisting}
*(Note : Le comportement de Mealy décrit ici assigne les sorties en fonction de l'état actuel et de l'entrée, bien que les valeurs de sortie soient similaires à celles de l'état actuel).*

\paragraph{4) Combien de bascules nécessite l'implantation de chaque machine ?}
Une bascule pour les deux.

\paragraph{On suppose maintenant que le signal t est asynchrone.}
1. Afin d'éviter l'apparition de métastabilité, décrire en VHDL un unique processus permettant de synchroniser le signal asynchrone \texttt{t}.

\begin{lstlisting}[language=VHDL]
-- Processus de synchronisation (Double bascule)
process(clk, reset)
begin
    if reset = '1' then
        t_sync1 <= '0';
        t_sync_out <= '0';
    elsif rising_edge(clk) then
        t_sync1 <= t;          -- Premiere bascule de synchronisation
        t_sync_out <= t_sync1; -- Deuxieme bascule de synchronisation
    end if;
end process;
\end{lstlisting}

\vspace{1cm}

\section*{Exercice 2 : Détecteur de séquence}

Dans cet exercice vous allez concevoir une machine à états qui détecte la séquence $abc= \text{<< } 111 \text{ >>}$ avec la contrainte que cette séquence doit être ordonnée temporellement : d'abord $a='1'$ pendant un premier cycle d'horloge, puis $b='1'$, puis enfin $c='1'$. La sortie $x$ devient '1' quand la séquence est correctement ordonnée. On considèrera les recouvrements (ex: $x$ reste à '1' tant que la séquence $111$ est maintenue).

\paragraph{1) Dessiner le diagramme de transitions de cette machine à états.}
La machine d'états (de type Moore) est conçue de manière à ce que les signaux $a, b, c$ apparaissent successivement. L'entrée est évaluée comme un vecteur $abc$.
\begin{itemize}
    \item \textbf{S0} (attente de $a$, $x=0$) : Si $abc = "100" \rightarrow \text{aller à S1}$. Sinon rester dans S0.
    \item \textbf{S1} (attente de $b$, $x=0$) : Si $abc = "110" \rightarrow \text{aller à S2}$. Si $abc = "100" \rightarrow \text{rester dans S1}$. Sinon $\rightarrow \text{retour à S0}$.
    \item \textbf{S2} (attente de $c$, $x=0$) : Si $abc = "111" \rightarrow \text{aller à S3}$. Si $abc = "110" \rightarrow \text{rester dans S2}$. Si $abc = "100" \rightarrow \text{aller à S1}$. Sinon $\rightarrow \text{retour à S0}$.
    \item \textbf{S3} (succès, $x=1$) : Si $abc = "111" \rightarrow \text{rester dans S3}$. Si $abc = "110" \rightarrow \text{aller à S2}$. Si $abc = "100" \rightarrow \text{aller à S1}$. Sinon $\rightarrow \text{retour à S0}$.
\end{itemize}
\textit{(Le schéma SVG correspondant est à inclure ici).}

\paragraph{2) Combien de bascules seront nécessaires au minimum pour implanter cette machine? Justifiez.}
Le diagramme de transitions comporte \textbf{4 états distincts} (S0, S1, S2, S3). Le nombre minimum de bascules $n$ nécessaires pour coder $E$ états est donné par la relation $2^n \ge E$. Ainsi, $2^n \ge 4$, ce qui donne \textbf{$n = 2$ bascules au minimum}.

\paragraph{3) Donner le code VHDL de cette machine.}

\begin{lstlisting}[language=VHDL]
library IEEE;
use IEEE.STD_LOGIC_1164.ALL;

entity sequence_detector is
    Port ( a   : in STD_LOGIC;
           b   : in STD_LOGIC;
           c   : in STD_LOGIC;
           clk : in STD_LOGIC;
           rst : in STD_LOGIC;
           x   : out STD_LOGIC);
end sequence_detector;

architecture Behavioral of sequence_detector is
    type state_type is (S0, S1, S2, S3);
    signal current_state, next_state : state_type;
    signal abc : std_logic_vector(2 downto 0);
begin
    -- Concaténation des entrees pour faciliter la lecture
    abc <= a & b & c;

    -- Processus synchrone
    process(clk, rst)
    begin
        if rst = '1' then
            current_state <= S0;
        elsif rising_edge(clk) then
            current_state <= next_state;
        end if;
    end process;

    -- Processus combinatoire
    process(current_state, abc)
    begin
        -- Valeurs par defaut pour eviter les latches
        next_state <= S0;
        x <= '0';

        case current_state is
            when S0 =>
                if abc = "100" then
                    next_state <= S1;
                end if;
            when S1 =>
                if abc = "110" then
                    next_state <= S2;
                elsif abc = "100" then
                    next_state <= S1;
                end if;
            when S2 =>
                if abc = "111" then
                    next_state <= S3;
                elsif abc = "110" then
                    next_state <= S2;
                elsif abc = "100" then
                    next_state <= S1;
                end if;
            when S3 =>
                x <= '1';
                if abc = "111" then
                    next_state <= S3;
                elsif abc = "110" then
                    next_state <= S2;
                elsif abc = "100" then
                    next_state <= S1;
                end if;
            when others =>
                next_state <= S0;
        end case;
    end process;
end Behavioral;
\end{lstlisting}

\paragraph{4) Combien de bascules générera la synthèse de votre code ? Justifiez.}
En utilisant un type énuméré VHDL (\texttt{type state\_type is (S0, S1, S2, S3);}), l'outil de synthèse (comme Vivado ou Quartus) choisit par défaut un encodage \textbf{"One-Hot"}. 
\begin{itemize}
    \item Si l'encodage \textbf{One-Hot} est appliqué, la synthèse générera \textbf{4 bascules} (une bascule dédiée pour chaque état).
    \item Si l'outil de synthèse est configuré pour optimiser la surface (encodage \textbf{Binaire}), il générera \textbf{2 bascules}, car 2 bits suffisent pour représenter 4 états.
\end{itemize}

\end{document}